\documentclass[10pt,a4paper]{amsart}
\usepackage[margin=19mm]{geometry}
\usepackage{amsmath,amssymb,mathtools}
\usepackage[scr=rsfs,cal=euler]{mathalfa}
\usepackage{xfrac}
\usepackage[unicode,hidelinks,bookmarks=false]{hyperref}
\newcommand{\ie}{i.e.\ }
\newcommand{\cf}{cf.\ }
\newcommand{\resp}{respectively\ }
\newcommand{\Subsection}{\subsection}
\providecommand{\nobreakdash}{\nobreak\hbox{-}}
\setlength{\emergencystretch}{2em}
\multlinegap0pt
\usepackage{amssymb}
\usepackage{esint}
\usepackage{enumerate}
\usepackage{mathtools}
\usepackage[shortlabels]{enumitem}
\usepackage{xcolor}
\usepackage{cancel}
\usepackage[normalem]{ulem}
\usepackage{caption}
\usepackage{tikz}
\newtheorem{lem}{Lemma}
\title{Partial regularity of the gradient for subsolutions}
\author{Aram Hakobyan, Michael Poghosyan, Henrik Shahgholian}
\date{Source: arXiv:2602.14305v2 (CC BY 4.0)}
\begin{document}
\maketitle

\noindent\emph{Source and licence.} Partial regularity of the gradient for subsolutions. Version arXiv:2602.14305v2 (CC BY 4.0), distributed by arXiv under the Creative Commons Attribution 4.0 International licence, http://creativecommons.org/licenses/by/4.0/, as recorded in the arXiv licence metadata for this exact version and shown on the abstract page https://arxiv.org/abs/2602.14305v2. Full licence notice: \texttt{licenses/arxiv-2602.14305-CC-BY-4.0.txt}.\par\smallskip

\noindent\textit{Editorial scope.} Counted unit: the article's lem lem:estimates (source0.tex lines 857--862). The article's complete original proof of it is delivered with it (source0.tex lines 864--892, one \texttt{proof} environment of 29 lines). No mathematics is written, repaired or re-derived: the statement and the proof are the article's own lines, in the article's own notation, and no retained line is substituted or re-flowed.  Retained uncounted as the source-local context the proof needs: lines 447--450; lines 472--474.  Nothing was omitted from any retained span, so the package carries no omission record.  No retained line contains a citation, so the wrapper carries no bibliography and no bibliography entry.  No retained line contains a figure environment, an image-inclusion command or drawing code, so no image is delivered and the package declares no asset dependency.  Editorial formatting only, in addition: an AMS \texttt{amsart} preamble that restates the article's own declarations and macros used by the retained lines, namely the environments \texttt{lem} on separate counters, together with the standard packages those lines need.  The retained environments print in this document as Lemma 1 in order of appearance, whereas the article prints the same items with the numbering of its own full document.  The complete native source of the article as distributed in the arXiv e-print for this exact version is bundled unchanged as \texttt{sources/194699/source0.tex}.
\begin{equation}\label{eq:distance}
 |y - y^0| < \epsilon^2, \qquad 
\max_{z\in \partial D_{y}^+}{\hbox{dist}(z, \partial D_{y^0}^+)} < \epsilon^2
\end{equation}

\begin{equation}\label{eq:estimate}
    I(r_\epsilon,y,u_y) =  I(r_\epsilon,y^0,u_y) + o_\epsilon (1).
\end{equation}

\begin{lem}\label{lem:estimates}
 We prove    \eqref{eq:estimate}, i.e.
 \begin{equation}
    I(r_\epsilon,y,u_y) =  I(r_\epsilon,y^0,u_y) + o_\epsilon (1).
\end{equation}
\end{lem}

\begin{proof}
This can be done by a direct calculation. We have 
$$
\Big| I(r_\epsilon,y,u_{y}) -  I(r_\epsilon,y^0,u_y)| \le
\frac{1}{r_\epsilon^2 }\int_{B_{r_\epsilon} (y) \setminus 
B_{r_\epsilon} (y^0)}  \frac{|\nabla u|^2}{|x-y|^{n-2}} + 
\frac{1}{r_\epsilon^2 }\int_{B_{r_\epsilon} (y^0) \setminus 
B_{r_\epsilon} (y)}  \frac{|\nabla u|^2}{|x-y^0|^{n-2}}
+
$$
$$
+\frac{1}{r_\epsilon^2 }\int_{B_{r_\epsilon} (y^0) \cap 
B_{r_\epsilon} (y)} |\nabla u|^2\left|\frac{1}{|x-y|^{n-2}}- \frac{1}{|x-y^0|^{n-2}}\right|.
$$
Since $|y -y^0 | < \epsilon^2$,  the  integration domains in 
the first two integrals on the right side of the inequality are small, and tend to zero as $\epsilon $ tends to zero.
 As to the last integral, it can be estimated in the following way, using the mean-value theorem and \eqref{eq:distance}:
$$
\frac{1}{r_\epsilon^2 }\int\limits_{B_{r_\epsilon} (y^0) \cap 
B_{r_\epsilon} (y)} |\nabla u|^2\left|\frac{1}{|x-y|^{n-2}}- \frac{1}{|x-y^0|^{n-2}}\right| \le C  \frac{|y-y^0|}{r_\epsilon^2 }\int\limits_{B_{r_\epsilon} (y^0) \cap 
B_{r_\epsilon} (y)} \frac{1}{\min\{|x-y|, |x-y^0|\}^{n-1}}\le 
$$
$$
\le 2C
\frac{|y-y^0|}{r_\epsilon^2 }\int_{B_{r_\epsilon} (y^0)} \frac{1}{|x-y^0|^{n-1}} =  C_1\frac{|y-y^0|}{r_\epsilon} = 
C_1\frac{|y-y^0|}{\epsilon+|y-y^0|}\le C_1\frac{|y-y^0|}{\epsilon}\le C_1\frac{\epsilon^2}{\epsilon} = C_1\epsilon. 
$$

\end{proof}

\end{document}
